% ============================================================
% Chapter 1: Introduction
% ============================================================
\chapter{Introduction}
\label{ch:introduction}

\section{Background and Motivation}
\label{sec:background}

The global imperative to transition toward sustainable energy systems has intensified research efforts across power systems engineering, building science, and transportation electrification. As energy infrastructure grows increasingly complex---with distributed generation, demand-responsive loads, and multi-scale coordination requirements---traditional model-based control approaches face fundamental limitations in scalability, adaptability, and computational tractability. Reinforcement learning (RL), with its capacity to learn complex control policies from interaction data without requiring explicit system models, has emerged as a promising paradigm for sequential decision-making in energy systems.

Recent years have witnessed a proliferation of RL applications in sustainable energy domains. In electric vehicle (EV) charging management, RL agents have demonstrated the ability to coordinate charging schedules across large networks while balancing grid constraints, user satisfaction, and carbon emissions~\citep{lee2019acndata, qiu2023marl_ev}. In building energy management, RL-based HVAC controllers have achieved significant energy savings while maintaining occupant comfort in commercial and residential buildings~\citep{wei2017deep_hvac, zhang2019model_hvac}. In power plant operations, RL has been applied to optimize the dispatch and coordination of gas turbines, steam turbines, and auxiliary equipment in combined heat and power (cogeneration) facilities~\citep{vu2021gas_turbine}.

Despite this growing body of work, a critical gap persists in the systematic evaluation of RL algorithms across sustainable energy domains. Most studies evaluate a single algorithm on a single environment under idealized conditions, providing limited insight into algorithm robustness, scalability, and comparative performance. Three fundamental challenges remain largely unaddressed:

\paragraph{Perturbation Robustness.} Real-world energy systems are inherently subject to perturbations. Sensors degrade, communication channels introduce latency and corruption, actuators exhibit imprecision, and environmental conditions fluctuate stochastically. Yet, RL agents are typically trained and evaluated in deterministic or low-perturbation simulation environments. The sensitivity of learned policies to state perturbation (sensor uncertainty), action perturbation (actuator imprecision), and dynamics perturbation (stochastic dynamics) remains poorly characterized. Understanding how different perturbation channels affect policy performance is essential for deploying RL controllers in physical energy infrastructure.

\paragraph{Safety Constraints.} Energy systems operate under strict physical, regulatory, and operational constraints. Voltage limits must be respected in power distribution networks, temperature bounds must be maintained in occupied buildings, and equipment operating ranges must not be violated in power plants. Standard RL algorithms, which optimize unconstrained cumulative reward, provide no mechanism for enforcing such constraints during training or deployment. Constrained Markov Decision Processes (CMDPs)~\citep{altman1999constrained} and the associated safe RL algorithms offer a principled framework for incorporating safety constraints, yet their application to sustainable energy domains remains limited.

\paragraph{Multi-Agent Coordination.} Many energy systems are naturally decomposable into multiple interacting subsystems---individual charging stations in an EV network, thermal zones in a multi-zone building, or gas turbines in a cogeneration plant. While single-agent RL treats the entire system as a monolithic control problem, multi-agent RL (MARL) enables distributed decision-making that can improve scalability, reduce communication overhead, and provide modularity. However, the comparative performance of single-agent versus multi-agent formulations in energy systems has not been systematically studied.

The SustainGym suite, introduced by \citet{yeh2023sustaingym}, provides a collection of RL environments grounded in real-world sustainable energy systems. SustainGym includes environments for EV charging (based on the Adaptive Charging Network at Caltech), building energy management (based on EnergyPlus RC thermal models with ASHRAE 90.1-2019 prototypes), and cogeneration plant control (based on ONNX surrogate models of industrial combined heat and power facilities). These environments offer realistic dynamics, physics-based models, and real-world data, making them suitable testbeds for rigorous RL benchmarking.

However, the original SustainGym work focused primarily on demonstrating environment functionality with basic RL baselines (PPO, SAC), without systematic investigation of perturbation robustness, safety constraints, or multi-agent coordination. This thesis addresses these gaps through a comprehensive benchmarking study that we term \textbf{SustainRL-Bench}.

\section{Problem Statement}
\label{sec:problem}

This thesis investigates the following overarching research question:

\begin{quote}
\textit{How do state-of-the-art reinforcement learning algorithms perform across diverse sustainable energy environments when subjected to realistic perturbations, safety constraints, and multi-agent decompositions?}
\end{quote}

Specifically, we consider three sustainable energy control problems, each formulated as a Markov Decision Process (MDP) or its extensions:

\begin{enumerate}
    \item \textbf{Electric Vehicle Charging Coordination}: Given a network of $n$ charging stations with time-varying arrivals, departures, energy demands, and carbon intensity signals, find a charging policy $\pi: \calS \to \calA$ that maximizes cumulative profit while minimizing carbon emissions and network constraint violations over a 24-hour horizon.

    \item \textbf{Building HVAC Control}: Given a multi-zone commercial building with coupled thermal dynamics governed by a resistance-capacitance (RC) model, find a control policy that minimizes the weighted sum of energy consumption and thermal comfort violations across all zones over a 24-hour horizon.

    \item \textbf{Cogeneration Plant Dispatch}: Given a combined heat and power facility with three gas turbines, a steam turbine, condenser, and cooling system, find a dispatch policy that minimizes fuel costs, ramping penalties, demand shortfalls, and constraint violations over a 24-hour operating horizon.
\end{enumerate}

Each problem is studied along three experimental axes:

\begin{description}
    \item[Axis~1 (Perturbation Robustness):] Systematic injection of state perturbation, action perturbation, and dynamics perturbation at varying intensities to characterize policy degradation curves and identify algorithm-specific failure modes.

    \item[Axis~2 (Safe RL):] Reformulation of each environment as a Constrained MDP with environment-specific cost signals, and evaluation of safe RL algorithms that satisfy constraint budgets while maintaining competitive reward performance.

    \item[Axis~3 (Multi-Agent RL):] Decomposition of each environment into a cooperative multi-agent system with per-component agents, and comparison of centralized single-agent versus decentralized multi-agent policy performance.
\end{description}

\section{Research Questions and Objectives}
\label{sec:research_questions}

The specific research questions addressed in this thesis are:

\begin{description}
    \item[RQ1 (Perturbation Sensitivity):] How sensitive are PPO, SAC, and TD3 to state, action, and dynamics perturbation across the three sustainable energy environments? Do different algorithms exhibit different perturbation robustness profiles, and do these profiles depend on the environment characteristics?

    \item[RQ2 (Perturbation Decomposition):] What is the relative contribution of each perturbation channel (state, action, dynamics) to policy degradation? Can we identify which perturbation channel is most critical for each environment-algorithm pair?

    \item[RQ3 (Safe RL):] Can safe RL algorithms (e.g., PPOLag, CPO, OnCRPO, FOCOPS) achieve competitive reward performance while satisfying environment-specific safety constraints? What is the Pareto frontier between reward and constraint satisfaction?

    \item[RQ4 (Multi-Agent RL):] Does multi-agent decomposition improve or degrade performance compared to centralized single-agent control? How does the choice of agent granularity (e.g., per-station, per-zone, per-turbine) affect learning efficiency and final performance?

    \item[RQ5 (Cross-Environment Insights):] Are there universal patterns in algorithm performance and perturbation sensitivity across the three environments, or is performance highly environment-specific?
\end{description}

The objectives of this thesis are:

\begin{enumerate}
    \item To develop a unified benchmarking framework (SustainRL-Bench) that enables systematic evaluation of RL algorithms across multiple sustainable energy environments.

    \item To implement and validate perturbation injection mechanisms for three independent perturbation channels (state, action, dynamics) in each environment.

    \item To formulate and implement CMDP wrappers for safe RL evaluation in all three environments.

    \item To implement multi-agent versions of all three environments using the PettingZoo parallel environment interface.

    \item To conduct a comprehensive experimental evaluation covering standard RL, safe RL, and MARL across all three environments with rigorous statistical methodology.

    \item To provide actionable insights for practitioners on algorithm selection, perturbation management, and multi-agent design in sustainable energy RL applications.
\end{enumerate}

\section{Contributions}
\label{sec:contributions}

The principal contributions of this thesis are as follows:

\begin{description}
    \item[C1. Unified Three-Axis Benchmarking Framework.] We present SustainRL-Bench, the first comprehensive RL benchmark that simultaneously evaluates perturbation robustness, safety constraints, and multi-agent coordination across multiple sustainable energy environments. Unlike prior work that addresses these axes in isolation, our framework enables cross-axis and cross-environment comparisons within a consistent experimental methodology.

    \item[C2. Systematic Perturbation Robustness Analysis.] We develop a three-channel perturbation injection framework that independently perturbs observations, actions, and environment dynamics. For each environment, we design physically motivated perturbation models: proportional Gaussian noise for sensor measurements, multiplicative noise for actuator commands, and stochastic perturbations of environmental parameters. We characterize policy degradation curves across 10+ perturbation levels for three algorithms (PPO, SAC, TD3), revealing algorithm-specific and environment-specific failure modes.

    \item[C3. Safe RL Formulation for Energy Systems.] We formulate each of the three environments as a CMDP by designing environment-specific cost signals that capture physically meaningful safety constraints: network capacity violations in EV charging, temperature comfort violations in buildings, and dynamic operational constraint violations in cogeneration. We evaluate safe RL algorithms from the OmniSafe framework, providing the first systematic comparison of constrained RL methods across sustainable energy domains.

    \item[C4. Multi-Agent Environment Design and Evaluation.] We design and implement multi-agent versions of all three environments using the PettingZoo parallel environment interface, with physically motivated agent decompositions: per-station agents in EV charging (54 agents), per-zone agents in building HVAC, and per-turbine agents in cogeneration (4 agents). We compare single-agent versus multi-agent performance using Ray RLlib, providing insights into the scalability-performance tradeoff.

    \item[C5. Reproducible Experimental Infrastructure.] We provide a complete, open-source experimental infrastructure including training scripts, evaluation protocols, SLURM job specifications for HPC deployment, and statistical analysis tools. All code and configurations are publicly available at \url{https://github.com/MehmetKORUTURK/sustaingym-rl-benchmark}. All experiments use standardized hyperparameters, fixed random seeds, and bootstrap confidence intervals following the recommendations of \citet{agarwal2021deep} for reliable RL evaluation.
\end{description}

\section{Thesis Organization}
\label{sec:organization}

The remainder of this thesis is organized as follows:

\textbf{Chapter~\ref{ch:literature}: Literature Review} provides a comprehensive review of related work across five domains: reinforcement learning for power and energy systems, perturbation robustness and distribution shift in RL, safe reinforcement learning, multi-agent reinforcement learning, and RL benchmarking methodologies. We position our contributions relative to the existing literature and identify the specific gaps that this thesis addresses.

\textbf{Chapter~\ref{ch:preliminaries}: Preliminaries and Background} establishes the mathematical foundations required for the subsequent technical chapters. We formally define Markov Decision Processes, derive the policy gradient theorem, and present the three deep RL algorithms used in this work (PPO, SAC, TD3). We then extend the MDP framework to Constrained MDPs for safe RL, introduce multi-agent extensions including stochastic games and decentralized partially observable MDPs, and formalize the perturbation models used throughout the thesis. Finally, we provide an overview of the SustainGym environment suite.

\textbf{Chapter~\ref{ch:methodology}: Methodology --- The SustainRL-Bench Framework} presents the core technical contribution of this thesis. We describe each of the three environments in detail, including their physical models, observation and action spaces, reward formulations, and constraint definitions. We then present the perturbation injection framework, safe RL formulations, and multi-agent decompositions for each environment. The chapter concludes with the experimental design, including training configurations, evaluation protocols, statistical methodology, and baseline methods.

\textbf{Chapter~\ref{ch:results}: Experimental Results and Discussion} presents the results of our comprehensive benchmarking study, organized by experimental axis: perturbation robustness results (including degradation curves and algorithm-specific failure modes), safe RL results (including reward--cost orthogonality analysis and negative findings), and MARL results (including single-agent versus multi-agent comparisons). Cross-environment patterns and actionable insights are discussed.

\textbf{Chapter~\ref{ch:conclusion}: Conclusion and Future Work} summarizes the key findings of this thesis, discusses limitations, and outlines directions for future research including transfer learning across environments, curriculum-based perturbation training, and real-world deployment considerations.
